\section{Hamiltonian-cycle structure}
\label{sec:hamiltonian-cycles}

The general classification below agrees with the complete certified
Hamiltonian-cycle universes of every analyzed double-ladder member, including
the prediction-locked graphs beyond order 36.  Those finite universes provide
an independent check of the cycle counts and classes proved here.

\begin{theorem}[Hamiltonian cycles of the double ladder]
\label{thm:cycle-formula}
For odd \(a,b\ge3\), the Hamiltonian cycles of \(D(a,b)\) are one all-rails
cycle, the \(ab\) cycles with one turn in each ladder, and four two-connector
exceptions.  In particular,
\[
 |\Hcal(D(a,b))|=ab+5.
\]
\end{theorem}

\begin{proof}
For a ladder of odd length \(m\), put binary variables \(x_{t,s}\) on rail
\(s\in\{0,1\}\) of cell \(t\), where \(0\le t<m\), and \(y_t\) on rung
\(t\), where \(0\le t\le m\).  Let \(c_{0,s},c_{m,s}\) be its endpoint
connector bits.  A Hamiltonian cycle satisfies the degree-two equations
\begin{align}
 x_{0,s}+y_0+c_{0,s}&=2,\label{eq:ladder-left-degree}\\
 x_{t-1,s}+x_{t,s}+y_t&=2 &&(0<t<m),\label{eq:ladder-inner-degree}\\
 x_{m-1,s}+y_m+c_{m,s}&=2.\label{eq:ladder-right-degree}
\end{align}
Writing \(d_t=x_{t,0}-x_{t,1}\) and subtracting the two side equations gives
\[
 d_0=c_{0,1}-c_{0,0},\qquad d_t=-d_{t-1},\qquad
 d_{m-1}=c_{m,1}-c_{m,0}.
\]
Since \(m-1\) is even, a necessary endpoint condition is
\begin{equation}
\label{eq:ladder-endpoint-compatibility}
 c_{0,0}-c_{0,1}=c_{m,0}-c_{m,1}.
\end{equation}
The connector bits therefore determine the rail differences.  A difference
of \(1\) or \(-1\) fixes both binary rail variables, whereas a difference
of zero leaves the choices \((0,0)\) and \((1,1)\).  Once the rail variables
are fixed, the degree equations determine every rung.  We impose connectivity
below to resolve the remaining choices.

The four connectors form the cut between the two nonempty ladders.  A
Hamiltonian cycle crosses this cut a positive even number of times, so it
uses either all four connectors or exactly two.

\smallskip\noindent\emph{Four connectors.}
Here all endpoint bits are one, so \(x_{t,0}=x_{t,1}=z_t\), and
\[
 y_0=1-z_0,\qquad y_m=1-z_{m-1},\qquad
 y_t=2-z_{t-1}-z_t\quad(0<t<m).
\]
Adjacent zero values of \(z\) would force a rung value of two and are
impossible.  If there are at least two zero positions, choose consecutive
ones \(i<j\) in their ordered list.  Both rails are present in cells
\(i+1,\ldots,j-1\), and rungs \(i+1\) and \(j\) are present.  These edges
form a cycle isolated from the rest of the graph by the omitted rail pairs
in cells \(i\) and \(j\).  Thus each ladder has either no zero position or
exactly one.

With no zero position, the ladder consists of its two straight rails.
With one zero position, it consists of two vertex-disjoint paths that together
span the ladder, each joining the two terminals at one end.  Number the connectors
\(1,2,3,4\) in the order of \cref{eq:connectors}.  The resulting terminal
pairings are
\[
\begin{array}{c|cc}
 &\text{straight}&\text{one turn}\\ \hline
 A&13\mid24&12\mid34\\
 B&12\mid34&13\mid24
\end{array}
\]
where, for example, \(13\mid24\) pairs terminals 1 with 3 and 2 with 4.
Different pairings on the two ladders join into a single spanning cycle;
identical pairings give two components.  Consequently both ladders must be
straight or both must turn.  The former gives one all-rails cycle \(R\).
The latter gives one cycle \(T(i,j)\) for each choice of turn cells
\(0\le i<a\) and \(0\le j<b\), hence exactly \(ab\) cycles.

\smallskip\noindent\emph{Two connectors.}
There are six connector words of weight two.  For \(1001\) and \(0110\),
the endpoint differences on ladder \(A\) have opposite signs, violating
\cref{eq:ladder-endpoint-compatibility}.  Thus these two words cannot even
satisfy the degree equations.  The four remaining words give the following
configurations:
\[
\begin{array}{c|cc}
\text{word}& A&B\\ \hline
0011&\text{all rails, rung }0&\text{alternating rails, all rungs}\\
1100&\text{all rails, rung }a&\text{alternating rails, all rungs}\\
0101&\text{alternating rails, all rungs}&\text{all rails, rung }0\\
1010&\text{alternating rails, all rungs}&\text{all rails, rung }b
\end{array}
\]
To justify uniqueness, in the alternating ladder the rail differences are
alternately \(1\) and \(-1\).  They fix each rail pair, and the degree
equations force every rung to be present.  These edges form a single spanning
path between the two selected connector terminals.

In the other ladder the two selected connectors are at the same end.  The
endpoint equations at the opposite, connector-free end force its rung and
both adjacent rails to be present.  All rail differences are zero.  If a
zero rail pair occurred, the portion from the connector-free end to the
first such pair would form an isolated cycle.  Connectivity therefore
forces all rails to be present; the degree equations leave only the rung
at the connector-free end.  This also gives a single spanning path between
the selected terminals.  The two connectors join the two paths into one
Hamiltonian cycle.  Each of the four words therefore has exactly one
Hamiltonian completion.

The cases exhaust all possible connector selections.  The turn positions
and connector words distinguish the resulting edge sets, giving precisely
\(1+ab+4=ab+5\) Hamiltonian cycles.
\end{proof}
